\section{Preliminaries on quandles}
\label{section: preliminaries on quandles}


In this section, we recall the basic definitions and fix the examples and notation used later. The reader may consult \cite{joyce_1982_a_classifying_invariant_of_knots_the_knot_quandle,
eisermann_2014_quandle_coverings_and_their_galois_correspondence,
preprint_valeriy_mohamed_mahender_2025_yangbaxter_equation_and_related_algebraic_structures} for more background.

\subsection{Quandles and fundamental notions}
\label{subsection: quandles and standard examples}

\begin{dfn}
\label[dfn]{definition: quandle}
    A \emph{quandle} $(Q,s)$ is a pair consisting of a set $Q$ and a map $s\colon Q\to\Map(Q,Q)$, $x\mapsto s_x$, satisfying the following conditions for all $x,y\in Q$:
\begin{enumerate}
    \item\label{condition:Q1} $s_x(x)=x$;
    \item\label{condition:Q2} $s_x$ is bijective;
    \item\label{condition:Q3} $s_x\circ s_y=s_{s_x(y)}\circ s_x$.
\end{enumerate}
\end{dfn}

\begin{dfn}
\label[dfn]{definition: quandle homomorphism}
    Let $(Q,s)$ and $(Q',s')$ be quandles. A map $f\colon Q\to Q'$ is a \emph{quandle homomorphism} if $f\circ s_x=s'_{f(x)}\circ f$ for every $x\in Q$. The category of quandles and quandle homomorphisms is denoted by $\Quandle$.
\end{dfn}

The condition~\eqref{condition:Q3} says each symmetry map $s_x$ is a quandle homomorphism $Q\to Q$.

\begin{eg}[Trivial quandles]
\label[eg]{example: trivial quandle}
    Every set $X$ has a quandle structure given by $s_x=\id_X$ for all $x\in X$. We call it the \emph{trivial quandle} structure on $X$.
\end{eg}

\begin{dfn}
\label[dfn]{definition: inner and transvection groups}
    Let $Q$ be a quandle.
\begin{enumerate}
    \item The group of quandle automorphisms of $Q$ is denoted by $\Aut(Q)$.
    \item The \emph{inner automorphism group} of $Q$ is $\Inn(Q)=\langle s_x\mid x\in Q\rangle\subseteq\Aut(Q)$.
    \item The \emph{transvection group} of $Q$ is $\Trans(Q) = \langle s_xs_y^{-1}\mid x,y\in Q\rangle$.
    \footnote{It is also called the \emph{displacement group} in some literature.}
    \item The quandle $Q$ is called \emph{connected} if the natural action of $\Inn(Q)$ on $Q$ is transitive. The $\Inn(Q)$-orbits are called the connected components of $Q$, and their set is denoted by $\pi_0(Q)$.
\end{enumerate}
\end{dfn}


For later quantitative statements, we recall finite generation. A subset $R\subseteq Q$ of a quandle $Q$ is called a \emph{subquandle} when $R$ is closed under both $s_z$ and $s_z^{-1}$ for every $z\in R$. A subset $S\subseteq Q$ \emph{generates} $Q$ if $Q$ is the smallest subquandle containing $S$. The quandle $Q$ is \emph{finitely generated} if it admits a finite generating set.

\begin{lem}
\label[lem]{lemma: quandle generators generate inner group}
    Let $Q$ be a quandle. If $Q$ is generated by $x_1,\ldots,x_r$ as a quandle, then $\Inn(Q)$ is generated by $s_{x_1},\ldots,s_{x_r}$.
\end{lem}

\begin{proof}
Let $H=\langle s_{x_1},\ldots,s_{x_r}\rangle \subseteq \Inn(Q)$. Since $s_{\varphi(x_i)}=\varphi s_{x_i}\varphi^{-1}\in H$ for $\varphi\in H$, the subset $H\cdot\{x_1,\ldots,x_r\}$ is closed under the quandle operation. So it is a subquandle of $Q$, and hence we have $Q=H\cdot\{x_1,\ldots,x_r\}$ by the assumption. Then, for every $x\in Q$, if we take an index $i$ and $\varphi\in H$ with $\varphi(x_i)=x$, it follows that $s_x=s_{\varphi(x_i)} \varphi s_{x_i} \varphi^{-1}\in H$, which shows that $\Inn(Q)=H$.
\end{proof}


We use the notation $\As(Q)$ for the adjoint group of a quandle. The following definition and properties are standard; see \cite[Sections~2 and 5]
{eisermann_2014_quandle_coverings_and_their_galois_correspondence}.

\begin{dfn}
\label[dfn]{definition: adjoint group}
    Let $Q$ be a quandle. The \emph{adjoint group} $\As(Q)$ is the group generated by $\{g_{x}\}_{x\in Q}$ subject to the relations $g_{x}g_{y}=g_{s_x(y)}g_{x}$, for $x,y\in Q$.
\end{dfn}

\begin{prop}
\label[prop]{proposition: elementary properties of adjoint groups}
Let $Q$ be a quandle.
\begin{enumerate}
    \item The assignment $g_{x}\mapsto s_x$ induces a surjective group homomorphism $\rho_{Q}\colon\As(Q)\twoheadrightarrow\Inn(Q)$, and $\As(Q)$ acts on $Q$ via $\rho_Q$.
    \item Every quandle homomorphism $f\colon P\to Q$ induces a group homomorphism $\As(f)\colon\As(P)\to\As(Q)$ satisfying $\As(f)(g_{x})=g_{f(x)}$. If $f$ is surjective, then so is $\As(f)$.
    \item The homomorphism $\rho_{Q}\colon\As(Q)\twoheadrightarrow\Inn(Q)$ is a central extension of groups.
\end{enumerate}
\end{prop}

There is a natural degree homomorphism $\varepsilon_{Q}\colon\As(Q)\to\mathbb Z$ defined by $\varepsilon_{Q}(g_{x})=1$; see \cite[Definition~2.33]{eisermann_2014_quandle_coverings_and_their_galois_correspondence}. We put $\As_0(Q)\coloneqq \Ker(\varepsilon_{Q})$. If $Q$ is connected, the group $\As_0(Q)$ agrees with the commutator subgroup of $\As(Q)$ \cite[Remark~2.34]{eisermann_2014_quandle_coverings_and_their_galois_correspondence}; this equality need not hold for a disconnected quandle.



\subsection{Surjective homomorphisms and relative symmetries}
\label{subsection: quandle homomorphisms}

In this subsection, we recall the results on surjective quandle homomorphisms used throughout the paper. Our main sources are Eisermann's covering theory \cite{eisermann_2014_quandle_coverings_and_their_galois_correspondence}, the quotient and factorization results of \cite{bunch_lofgren_rapp_yetter_2010_on_quotients_of_quandles,even_gran_2014_on_factorization_systems_for_surjective_quandle_homomorphisms}, and the relative symmetry groups introduced in \cite{preprint_yuki_tomoki_2026_relativization_of_symmetries_on_quandles}.

We shall use quotients by quandle congruences and by group actions.

\begin{dfn}
\label[dfn]{definition: quandle congruence}
An equivalence relation $\qCong$ on a quandle $Q$ is a \emph{quandle
congruence} if
\[
    x\mathrel\qCong x',\quad y\mathrel\qCong y'
    \quad\Longrightarrow\quad
    s_x^{\varepsilon}(y)\mathrel\qCong
    s_{x'}^{\varepsilon}(y')
\]
for $\varepsilon\in\{1,-1\}$.
\end{dfn}

\begin{prop}
%[{\cite[Propositions~1.22 and 1.23]{preprint_yuki_tomoki_2026_relativization_of_symmetries_on_quandles}}]
\label[prop]{proposition: quotients by quandle congruences}
    Let $\qCong$ be a quandle congruence on $Q$. Then, the quotient set $Q/{\qCong}$ has a unique quandle structure satisfying
    \[ s_{[x]}([y])=[s_x(y)], \]
    and the quotient map $Q\twoheadrightarrow Q/{\qCong}$ is a surjective quandle homomorphism. If $f\colon Q\twoheadrightarrow R$ is surjective, then its kernel relation
    \[ \Eq(f) = \{(x,y)\in Q\times Q\mid f(x)=f(y)\} \]
    is a quandle congruence. These correspondences are inverse to each other.
\end{prop}

\begin{prop}[{see \cite[Proposition~1.25 and Remark~1.26]{preprint_yuki_tomoki_2026_relativization_of_symmetries_on_quandles}}]
%\cite[Lemma~1.20]{darne_2026_nilpotent_quandles}; 
\label[prop]{proposition: orbit quotients of quandles}
    Let $A\subseteq\Aut(Q)$ be a subgroup normalized by $\Inn(Q)$. Then the relation given by the $A$-orbits is a quandle congruence. We denote the quotient quandle by $Q/A$. In particular, this construction applies to every normal subgroup $A\trianglelefteq\Inn(Q)$.
\end{prop}

The assignment $Q\mapsto\Inn(Q)$ is not functorial for arbitrary quandle homomorphisms, but it is functorial for surjections.

\begin{prop}[{\cite[Proposition~2.26]{eisermann_2014_quandle_coverings_and_their_galois_correspondence}}]
\label[prop]{proposition: functoriality of Inn for surjections}
    Let $f\colon P\twoheadrightarrow Q$ be a surjective quandle homomorphism. There is a unique surjective group homomorphism
    \[ f_*\colon\Inn(P)\twoheadrightarrow\Inn(Q) \]
    satisfying $f_*(s_x)=s_{f(x)}$ for every $x\in P$.
    Moreover, this gives rise to a functor $\Inn(-) \colon \Quandle^\mathrm{surj} \to \Grp^\mathrm{surj}$ from the category of surjective quandle homomorphisms to the category of surjective group homomorphisms.
\end{prop}

We next recall coverings in the sense of Eisermann.

\begin{dfn}[{\cite[Definition~2.42]{eisermann_2014_quandle_coverings_and_their_galois_correspondence}}]
\label[dfn]{definition: quandle covering}
    A quandle homomorphism $f\colon P\to Q$ is a \emph{covering homomorphism} if it is surjective and $f(x)=f(y)$ implies $s_x=s_y$ for all $x,y\in P$.
\end{dfn}

\begin{prop}[{\cite[Proposition~4.19]{eisermann_2014_quandle_coverings_and_their_galois_correspondence}}]
\label[prop]{proposition: pull-back of covering is again covering}
    Covering homomorphisms are preserved by pullback along arbitrary quandle homomorphisms.
\end{prop}

\begin{prop}[{\cite[Proposition~2.49]{eisermann_2014_quandle_coverings_and_their_galois_correspondence}}]
\label[prop]{proposition: covering hom is central extension}
    If $f\colon P\twoheadrightarrow Q$ is a covering homomorphism, then $f_*\colon\Inn(P)\twoheadrightarrow\Inn(Q)$ is a central extension of groups.
\end{prop}

We now recall the relative symmetry groups.

\begin{dfn}[{\cite[Definitions~2.1 and 3.1]{preprint_yuki_tomoki_2026_relativization_of_symmetries_on_quandles}}]
\label[dfn]{definition: relative inner and transvection groups}
    Let $f\colon P\twoheadrightarrow Q$ be a surjective quandle homomorphism.
\begin{enumerate}
    \item The \emph{relative inner automorphism group} of $f$ is
    \[ \relInn(f) = \Ker\bigl( f_*\colon \Inn(P)\to\Inn(Q) \bigr). \]
    \item The \emph{relative transvection group} of $f$ is
    \[ \relTrans(f) = \left\langle s_xs_y^{-1} \mid x,y\in P \text{ such that } f(x)=f(y) \right\rangle. \]
\end{enumerate}
\end{dfn}

\begin{prop}[{\cite[Propositions~2.4 and~3.4, and Lemma~3.15]{preprint_yuki_tomoki_2026_relativization_of_symmetries_on_quandles}}]
\label[prop]{lemma: elementary properties of relative symmetry groups}
    For every surjective homomorphism $f\colon P\twoheadrightarrow Q$, the following hold.
    \begin{enumerate}
        \item The group $\relInn(f)$ is normal in $\Inn(P)$ and $\relInn(f)=\{\varphi\in \Inn(P) \mid f\circ \varphi=f \}$. Every inner automorphism $\varphi\in \relInn(f)$ preserves the fibers of $f$, and hence $\relInn(f)$ acts on each fiber.
        \item The group $\relTrans(f)$ is normal in $\Inn(P)$ and $\relTrans(f)\subseteq\relInn(f)$. Moreover, $f$ is covering if and only if $\relTrans(f)=1$.
    \end{enumerate}
\end{prop}

\begin{rem}
\label[rem]{remark: relative transvection as relative displacement group}
    The group $\relTrans(f)$ is described in \cite{bonatto_stanovsky_2021_commutator_theory_for_racks_and_quandles} as the relative displacement group $\Dis_\alpha$ where $\alpha$ is the congruence corresponding to the surjection $f$.
\end{rem}

\begin{prop}[{\cite[Proposition~2.15]{preprint_yuki_tomoki_2026_relativization_of_symmetries_on_quandles}}]
\label{prop:basic_property_of_quotient_by_normal_subgroup}
Let $Q$ be a quandle and let $N\subseteq \Inn(Q)$ be a normal subgroup.
\begin{enumerate}
    \item\label{item:N_is_contained_in_relInn_of_projection}
    For the quotient map $\pi_N\colon Q\twoheadrightarrow Q/N$, we have $N \subseteq \relInn(\pi_N)=\Ker(\Inn(\pi_N))$.
    \item\label{item:factoring_condition_through_quotient_by_normal_subgroup}
    A quandle homomorphism $f\colon Q\to R$ factors through $\pi_N\colon Q\twoheadrightarrow Q/N$ if and only if $f\circ \varphi = f$ holds for all $\varphi\in N$, or equivalently if and only if $N\subseteq \relInn(f)$.
\end{enumerate}
\end{prop}

The following lemma will be used repeatedly.

\begin{lem}[{\cite[Lemma 3.8]{preprint_yuki_tomoki_2026_relativization_of_symmetries_on_quandles}}]
\label[lem]{lemma: relative symmetry groups under a surjective factorization}
    Suppose that $f=h\circ u$, where $u\colon P\twoheadrightarrow R$ and $h\colon R\twoheadrightarrow Q$ are surjective. Then
    \[ u_*(\relInn(f))=\relInn(h), \qquad u_*(\relTrans(f))=\relTrans(h).\]
\end{lem}


\begin{dfn}
\label[dfn]{definition: connected and rigid quandle homomorphisms}
    Let $f\colon P\twoheadrightarrow Q$ be a surjective quandle homomorphism.
\begin{enumerate}
    \item $f$ is called \emph{connected} if
    $\relInn(f)$ acts transitively on every fiber of $f$ (\cite[Definition~2.5]{preprint_yuki_tomoki_2026_relativization_of_symmetries_on_quandles}).
    \item $f$ is called \emph{strongly connected} if $\relTrans(f)$ acts transitively on every fiber of $f$ (\cite[Definition~3.6]{preprint_yuki_tomoki_2026_relativization_of_symmetries_on_quandles}).
    \item The homomorphism $f$ is called \emph{rigid} if $f_*\colon\Inn(P)\to\Inn(Q)$ is an isomorphism, or equivalently, if $\relInn(f)=1$ (\cite[Definition~4.1]{bunch_lofgren_rapp_yetter_2010_on_quotients_of_quandles}).
\end{enumerate}
\end{dfn}

We can see that a rigid homomorphism is covering, and that a connected rigid homomorphism is an isomorphism \cite[Lemmas~2.19 and~2.20]{preprint_yuki_tomoki_2026_relativization_of_symmetries_on_quandles}.
%Connectedness descends along surjective factorizations \cite[Proposition~2.12(3) and Remark~2.13]{preprint_yuki_tomoki_2026_relativization_of_symmetries_on_quandles}.

The following factorization is essentially \cite[Theorem~8.1]{bunch_lofgren_rapp_yetter_2010_on_quotients_of_quandles}, and its factorization-system formulation is due to \cite[Propositions~3.1 and~3.2]{even_gran_2014_on_factorization_systems_for_surjective_quandle_homomorphisms}.

\begin{thm}[Quandle Stein factorization; {\cite[Theorem~2.22 and Proposition~2.24]{preprint_yuki_tomoki_2026_relativization_of_symmetries_on_quandles}}]
\label[thm]{theorem: connected-rigid factorization}
    Let $f\colon P\twoheadrightarrow Q$ be a surjective quandle homomorphism. Then, the canonical factorization $P\twoheadrightarrow P/\relInn(f)\twoheadrightarrow Q$ has connected first homomorphism and rigid second homomorphism. 
    %
    Moreover, the classes of connected and rigid homomorphisms form an orthogonal factorization system for the surjections.
    % Remark~2.25: satisfying the Frobenius condition.
\end{thm}

In particular, rigid homomorphisms are preserved by pullback along surjective homomorphisms \cite[Proposition~2.21(2)]{preprint_yuki_tomoki_2026_relativization_of_symmetries_on_quandles}.

The relative transvection group gives the maximal covering factorization.

\begin{thm}[{
\cite[Propositions~3.17 and~3.18]{preprint_yuki_tomoki_2026_relativization_of_symmetries_on_quandles}}]
\label[thm]{theorem: maximal covering factorization}
    Let $f\colon P\twoheadrightarrow Q$ be a surjective quandle homomorphism. Then, the canonical factorization $P\twoheadrightarrow P/\relTrans(f)\twoheadrightarrow Q$ has connected first homomorphism and covering second homomorphism. It is universal among factorizations of $f$ whose second homomorphism is covering.
\end{thm}



Recall that a normal subgroup $K\trianglelefteq G$ is called \emph{$G$-perfect} if $[G,K]=K$. 

\begin{thm}[{\cite[Theorem 3.13 and Corollary 3.14]{preprint_yuki_tomoki_2026_relativization_of_symmetries_on_quandles}}]
\label{theorem: characterization of strongly connected via perfect subgroup}
    A surjective quandle homomorphism $f\colon P\twoheadrightarrow Q$ is strongly connected if and only if it is isomorphic to the quotient $\pi_N\colon P\twoheadrightarrow P/N$ by an $\Inn(P)$-perfect subgroup $N\subseteq \Inn(P)$.

    Moreover, there is a one-to-one correspondence between the isomorphism classes of strongly connected quotients of $P$ and $\Inn(P)$-perfect normal subgroups of $\Inn(P)$.
\end{thm}

\begin{cor}
\label{proposition:relTrans_of_strongly_conncted_is_Inn-perfect}
    For a strongly connected homomorphism $f\colon P\twoheadrightarrow Q$, the group $\relTrans(f)$ is $\Inn(P)$-perfect.
\end{cor}



We next recall the standard examples used in this paper.

\begin{eg}[Conjugacy quandles]
\label[eg]{example: conjugacy quandle}
    Let $G$ be a group. Its \emph{conjugacy quandle} $\Conj(G)$ is the set $G$ equipped with $s_x(y)=xyx^{-1}$ for $x,y\in G$. A surjective group homomorphism $\phi$ induces a surjective quandle homomorphism $f\coloneqq \Conj(\phi)\colon\Conj(G)\to\Conj(H)$. Write $K\coloneqq \Ker(\phi)$ and $A\coloneqq \phi^{-1}(Z(H))$. Then, we have
    \begin{align*}
        \relInn(f) &= \{ i_g \mid g \in G \text{ such that } \phi(g) \in Z(H)\} \cong A/Z(G), \text{ and} \\
        \relTrans(f) &= \{i_k \mid k\in \Ker(\phi)\} \cong K /(K\cap Z(G)),
    \end{align*}
    where $i_g$ is given by $i_g(x)=gxg^{-1}$ (see {\cite[Proposition~5.1]{preprint_yuki_tomoki_2026_relativization_of_symmetries_on_quandles}}).
\end{eg}


\begin{eg}[Alexander quandles]
\label[eg]{example: Alexander quandle}
    Let $A$ be an abelian group and let $\sigma\in\Aut(A)$ be a group automorphism of $A$. The \emph{Alexander quandle} $\Alex(A,\sigma)$ is the set $A$ equipped with
    \[ s_x(y)=(1-\sigma)x+\sigma(y) \qquad (x,y\in A). \]
    For $a\in A$, let $L_a\colon A\to A$ be the translation $L_a(y)=a+y$. It is immediate to see that $s_0=\sigma$ and for any $x\in \Alex(A,\sigma)$, $s_x = L_{(1-\sigma)x} \circ \sigma$ as maps.
    
    Let $(B,\tau)$ be another abelian group with an automorphism, and let $\Phi\colon A\twoheadrightarrow B$ be a surjective group homomorphism satisfying $\Phi\sigma=\tau\Phi$. It induces a surjective quandle homomorphism $f\coloneqq \Alex(\Phi)\colon \Alex(A,\sigma)\twoheadrightarrow\Alex(B,\tau)$.
    Write $K\coloneqq\Ker(\Phi)$ and $D\coloneqq(1-\sigma)A\cap K$. Then, we have
    \begin{align*}
        \relInn(f) &= \{L_a\sigma^n \mid a\in (1-\sigma)A\cap K,\ n\in\ZZ \text{ such that }\tau^n=\id_B\}, \text{ and} \\
        \relTrans(f) &= \{L_a \mid a\in(1-\sigma)K\} \cong (1-\sigma)K
    \end{align*}
    (see {\cite[Propositions~5.12 and 5.14]{preprint_yuki_tomoki_2026_relativization_of_symmetries_on_quandles}}).
    
    In particular, if the induced homomorphism $\langle\sigma\rangle \twoheadrightarrow \langle\tau\rangle$ is injective, then $\relInn(f)=\{L_a\mid a\in D\}$.
\end{eg}

\begin{eg}[Linear Alexander quandles]
\label[eg]{example: linear Alexander quandle}
    Let $n\geq1$, and let $a$ be an integer coprime to $n$. The Alexander quandle associated with multiplication by $a$ on $\ZZ/n\ZZ$ is denoted by $\Lambda_{n,a}$. Explicitly, the symmetry maps on the Alexander quandle $\Lambda_{n,a}$ are given by
    \[ s_x(y)=(1-a)x+ay \qquad (x,y\in\mathbb Z/n\mathbb Z). \]
    In particular, $\Lambda_{n,1}$ is the trivial quandle of order $n$ and $R_n\coloneqq \Lambda_{n, -1}$ is the dihedral quandle of order $n$.
\end{eg}
